% Generated by task tex-groundtruth from dossiers/gt-bus.md (sections A.6 and G15, summary item 1),
% with the corrections of dossiers/verify-gt-bus.md (its section 4, item 5, adds the last row).
\subsection{The setup, the commitment and the bus phase: disagreements between the three sources}\label{sec:gen-bus-disagreements}

\begin{sloppypar}
Short forms in this section. Specification files are under \file{doc/leanvm/body/}:
\code{05:} is \file{05-arithmetization.tex}, \code{06:} is \file{06-bus-interactions.tex},
\code{08:} is \file{08-end-to-end-protocol.tex}. Rust files are under \file{crates/lean_vm/src/}:
\code{mod.rs} is \file{cpu/mod.rs}, \code{layout.rs} is \file{cpu/layout.rs}, and \code{leaf.rs},
\code{constraints.rs} are at that root; \code{fs/} is \file{crates/fiat_shamir/src/}. \code{py:} is
\file{python-verifier/verifier.py}. The number after the colon is the line, at leanVM's pin
\code{a386121f}.
\end{sloppypar}

The three sources agree on the transcript of the setup, the commitment and the bus phase, up to
omissions of the specification (it does not mention the announced rate, the \code{BLAKE2S} floor,
the stacking window, the canonical encodings, the last unused combiner) and the prose order of
the two roots. The Rust and the Python verifier accept the same transcripts in these phases; no
divergence between them was found. The specification is silent on seven points the verifiers
fix, and its prose names the two roots in the other order. The blueprint follows the Rust on
each, rightly (\cref{f:bus-disagreements}).

{\small\setlength{\tabcolsep}{4pt}
\begin{longtable}{@{}>{\raggedright\arraybackslash}p{0.04\linewidth}>{\raggedright\arraybackslash}p{0.2\linewidth}>{\raggedright\arraybackslash}p{0.71\linewidth}@{}}
\caption{Disagreements between the specification, the Rust and the Python in the setup, the
commitment and the bus phase.}\label{tab:gen-bus-disagreements}\\
\toprule
\# & Subject & Disagreement \\
\midrule
\endfirsthead
\toprule
\# & Subject & Disagreement \\
\midrule
\endhead
\bottomrule
\endfoot
1 & The announced rate
  & The announced rate $ρ$ is sent with the sizes (\code{mod.rs:123, 149}; \code{py:1376});
  \code{08:55} names seven announced values. \\
2 & The order of the roots
  & \code{08:68} names $R_c$ first; the stream has $R$ first. \\
3 & The seed
  & \code{08:55} has ``Flock's BLAKE2s R1CS matrices''; the Rust hashes one constant naming the
  circuit (\code{mod.rs:66-93}). (The status file's finding on the seed, F14.) \\
4 & Checks of the verifiers the specification does not mention
  & The upper limbs of the announced sizes (check 2), the \code{BLAKE2S} floor (check 7), the rate window
  (check 8), the stacking window (check 9), the zero third limb of the root's halves (check 10), the length of
  $ζ$ (check 16), the consumption of the stream (check 17): the checks of \cref{tab:gen-bus-checks}. \\
5 & The combiner after the last layer
  & The verifiers draw one combiner more than layers (\cref{f:bus-extra-combiner}); \code{05:91}
  has one per layer. \\
6 & The program in the Python verifier
  & The Python verifier takes the program as the stacked table of $16·2^{κ_{\mathrm{bc}}}$ words
  and evaluates it whole (\code{py:566, 857}); it does not check that slots 0 to 2 and 11 to 15
  are zero. The Rust builds the eight columns from the instructions (\code{layout.rs:229-309}).
  The two agree on every table that is a program's; the Python's statement space is larger. Not
  a soundness matter: the table is the public statement. \\
7 & The fingerprint bound, the count columns, the assertions
  & The Rust's bound for the fingerprint is $5·2^μ$ (\code{leaf.rs:110}), the specification's
  $4·2^μ$ (\code{05:37}); and ``every count column'' (\code{06:78}) against the tables' read
  counts. The Rust's assertions on the layout (check 11 of \cref{tab:gen-bus-checks}) have no counterpart in the Python. \\
8 & The table sumcheck's round message, in the Rust's own comment (added by the citation check)
  & The comment at \code{constraints.rs:265-266} says ``\code{h(0)} is derived from the running
  claim rather than transmitted'', but the call at \code{:267},
  \code{next_round_poly(4, claim, None)}, transmits $c_0 = h(0)$ and derives $c_1$
  (\code{fs/transcript.rs:291, 299}). The comment is stale; the code and the Python agree with
  each other. \\
\end{longtable}
}
